\documentclass[a4paper,11pt]{article}
\usepackage[T2A]{fontenc}
\usepackage[utf8]{inputenc}
\usepackage[russian]{babel}
\usepackage{amsmath,amssymb,mathtools}
\usepackage{geometry}
\usepackage{microtype}
\usepackage{helvet}
\renewcommand{\familydefault}{\sfdefault}
\usepackage{tabularx,booktabs,array}
\usepackage{tikz}
\usepackage{xcolor}
\usepackage{hyperref}
\usepackage{fancyhdr}
\usepackage{parskip}
\usepackage{enumitem}
\geometry{left=23mm,right=23mm,top=23mm,bottom=25mm}
\definecolor{accent}{HTML}{176D75}
\definecolor{darkaccent}{HTML}{164D55}
\definecolor{textgray}{HTML}{3B4650}
\definecolor{errorred}{HTML}{A13C3C}
\hypersetup{colorlinks=true,linkcolor=accent,pdfauthor={Лёвин Артём Александрович},pdftitle={Ревью домашней работы — Даниил Хрястов — 05.11.2026}}
\setlength{\parindent}{0pt}
\setlength{\parskip}{0.65em}
\setlist[itemize]{leftmargin=1.5em,itemsep=0.2em,topsep=0.2em}
\pagestyle{fancy}
\fancyhf{}
\fancyhead[L]{\footnotesize\textcolor{textgray}{Математика \quad / \quad Ревью домашней работы}}
\fancyhead[R]{\footnotesize\textcolor{textgray}{05.11.2026}}
\fancyfoot[C]{\footnotesize\textcolor{textgray}{\thepage}}
\renewcommand{\headrulewidth}{0.25pt}
\renewcommand{\footrulewidth}{0pt}
\newcommand{\Err}{\textcolor{errorred}{\textbf{Ошибка.}}}
\newcommand{\Key}{\textcolor{darkaccent}{\textbf{Ключевая идея.}}}
\newcommand{\Result}{\textcolor{accent}{\textbf{Итог.}}}
\newcommand{\taskhead}[2]{\phantomsection\label{task:#1}\section*{Задание № #1\\[3pt]\large\normalfont #2}}
\begin{document}
\thispagestyle{empty}
\vspace*{37mm}
{\Large\textcolor{accent}{МАТЕМАТИКА}}\par
\vspace{13mm}
{\fontsize{30}{36}\selectfont\bfseries Ревью\\ домашней работы}\par
\vspace{7mm}
{\Large Степени, преобразования\\ и формулы с физическими величинами}\par
\vspace{24mm}
{\large Даниил Хрястов}\par
\vspace{5mm}
{\normalsize 5 ноября 2026 г.}\par
\vfill
{\normalsize Лёвин Артём Александрович}\par
{\small\textcolor{textgray}{эксклюзивно для Даниила Хрястова}}\par
\vspace{15mm}
\begin{center}
\begin{tikzpicture}[x=1cm,y=1cm]
\draw[accent!65,line width=1.0pt] (-5,0) .. controls (-2.7,0.7) and (-1.7,-0.7) .. (0,0)
.. controls (1.9,0.65) and (3.4,-0.48) .. (5,0.16);
\end{tikzpicture}
\end{center}
\newpage
\section*{Карта разбора}
По фотографиям проверены десять заданий. Ниже представлены только те задания, в которых обнаружены ошибки или существенные пробелы в решении. В №\,10 ответ получился верным, однако использованный переход требует исправления.\par
\vspace{3mm}
\renewcommand{\arraystretch}{1.3}
\small
\begin{tabularx}{\linewidth}{@{}p{13mm}p{29mm}X p{23mm}p{20mm}@{}}
\toprule
\textbf{№} & \textbf{Тема} & \textbf{Что исправить} & \textbf{Ваш ответ} & \textbf{Верный}\\
\midrule
\hyperref[task:2]{2} & Стандартный вид числа & Умножение степеней и коэффициентов & Не указан & $1000$\\
\hyperref[task:4]{4} & Формула и перенос & Знак при переносе $-mg$ & $0{,}6$ (с вопросом) & $0{,}6$ м\\
\hyperref[task:6]{6} & Степени десяти & Потеря порядка величины & $2$ & $2000$\\
\hyperref[task:10]{10} & Дробная степень & Неверное возведение дроби в степень & $0{,}002$ & $0{,}002$ м$^3$\\
\bottomrule
\end{tabularx}\par
\normalsize
\vspace{5mm}
\section*{Ответы на Ваши вопросы}
\subsection*{О задании № 2: «Куда переносить что?»}
Здесь \emph{ничего переносить между частями уравнения не нужно}: требуется умножить два числа в стандартном виде. Удобно выполнять действия параллельно: сначала обычные множители, затем степени десятки.
\[
(4\cdot10^{-7})(2{,}5\cdot10^9)
=\underbrace{(4\cdot2{,}5)}_{10}
\cdot\underbrace{10^{-7+9}}_{10^2}
=10^3=1000.
\]
Показатели $-7$ и $9$ складываются: $-7+9=2$. Само число $10^{-7}$ означает $1/10^7$.

\subsection*{О задании № 3: «Минус пропадает, потому что величина не может быть отрицательной?»}
В условии записано $x^{3/2}=125$ и отдельно указано $x>0$. Отрицательного показателя или минуса перед $x$ в исходной записи нет. Для положительного $x$ справедливо:
\[
 x^{3/2}=(\sqrt{x})^3=125=5^3
 \quad\Longrightarrow\quad\sqrt{x}=5\quad\Longrightarrow\quad x=25.
\]
\textbf{Минус нельзя просто удалить.} Если значение противоречит условию или области определения, его обоснованно исключают. Отрицательные числа вообще существуют и часто являются допустимыми ответами в других задачах. В данном уравнении решение $x=25$ удовлетворяет условию и даёт $25^{3/2}=125$.
\newpage
\taskhead{2}{Умножение чисел в стандартном виде}
\textbf{Условие.} Вычислите $(4\cdot10^{-7})\cdot(2{,}5\cdot10^9)$.

\textbf{Ваше решение.} В начале верно записан переход $2{,}5\cdot10^9=25\cdot10^8$. Далее заметен переход к сложению отдельных произведений; нижние вычисления зачёркнуты, поэтому надёжного итогового ответа нет.

\textbf{Ваш ответ.} Не указан в окончательном читаемом виде.

\Err{} При умножении двух произведений начали складывать части выражения. Результат зависит от того, \emph{все ли} множители были перемножены.

\textbf{Где именно.} Последний уверенно правильный шаг:
\[
(4\cdot10^{-7})(2{,}5\cdot10^9)
=(4\cdot10^{-7})(25\cdot10^8).
\]
Первый проблемный переход --- последующее разбиение произведения на сумму. Здесь требуется сохранить знак умножения между коэффициентами и степенями.

\textbf{Почему.} Если $a,b,c,d$ --- числа, то $(ab)(cd)=abcd$. Складывать разрешается показатели \emph{при умножении степеней с одним основанием}: $10^m\cdot10^n=10^{m+n}$. Коэффициенты $4$ и $2{,}5$ при этом умножаются между собой.

\Key{} Собрать вместе обычные числа, отдельно --- степени десятки.

\textbf{Исправление.} После корректной замены $2{,}5\cdot10^9=25\cdot10^8$ получаем
\[
(4\cdot25)\cdot10^{-7+8}=100\cdot10^1=1000.
\]

\textbf{Полное решение (короткий путь).}
\begin{align*}
(4\cdot10^{-7})\cdot(2{,}5\cdot10^9)
 &= (4\cdot2{,}5)\cdot(10^{-7}\cdot10^9)\\
 &= 10\cdot10^{-7+9}\\
 &= 10\cdot10^2=1000.
\end{align*}
\textbf{Проверка.} $4\cdot10^{-7}=0{,}0000004$, $2{,}5\cdot10^9=2\,500\,000\,000$; их произведение равно $1000$.

\Result{} $\boxed{1000}$.

\subsection*{Задача на закрепление}
Вычислите $(3\cdot10^{-4})\cdot(2\cdot10^7)$.
\newpage
\taskhead{4}{Выражение неизвестной из физической формулы}
\textbf{Условие.} $F=\rho g\ell^3-mg$, где $\rho=1000$ кг/м$^3$, $g=10$ м/с$^2$, $m=16$ кг, $F=2000$ Н. Найдите $\ell$ в метрах.

\textbf{Ваше решение.} Правильно получено $2000=10000\ell^3-160$, после чего записано $10000\ell^3=-2160$. Ниже появляются отрицательное значение $\ell^3$ и итог $0{,}6$ со знаком вопроса.

\textbf{Ваш ответ.} $0{,}6$ (записан с вопросом; ход решения противоречит ответу).

\Err{} Неверный знак после переноса числа $-160$ в другую часть равенства.

\textbf{Где именно.} Последний правильный шаг:
\[
2000=10000\ell^3-160.
\]
Первый неверный переход:
\[
10000\ell^3=\textcolor{errorred}{-2160}.
\]

\textbf{Почему.} Для сохранения равенства необходимо \emph{прибавить 160 к обеим частям}. Тогда справа $-160+160=0$, а слева $2000+160=2160$. У числа меняется знак как результат конкретной операции над обеими частями уравнения.

\Key{} Вместо запоминания направления «переноса» каждый раз выполняйте одинаковое действие с обеими частями равенства.

\textbf{Исправление и полное решение.}
\begin{align*}
2000&=1000\cdot10\cdot\ell^3-16\cdot10,\\
2000&=10000\ell^3-160,\\
2000+160&=10000\ell^3,\\
2160&=10000\ell^3,\\
\ell^3&=\frac{2160}{10000}=0{,}216
=\frac{27}{125}=\left(\frac35\right)^3.
\end{align*}
Длина положительна, поэтому $\ell=\frac35=0{,}6$ м.

\textbf{Проверка.}
\[
1000\cdot10\cdot(0{,}6)^3-16\cdot10
=2160-160=2000\ \text{Н}.
\]
\Result{} $\boxed{\ell=0{,}6\ \text{м}}$.

\subsection*{Задача на закрепление}
В формуле $F=\rho g\ell^3-mg$ известны $\rho=1000$ кг/м$^3$, $g=10$ м/с$^2$, $m=20$ кг и $F=7090$ Н. Найдите $\ell>0$.
\newpage
\taskhead{6}{Работа со степенями десяти при выражении температуры}
\textbf{Условие.} Из $P=\sigma S T^4$ найдите $T>0$, если $P=8\cdot10^6$, $\sigma=5\cdot10^{-8}$, $S=10$.

\textbf{Ваше решение.} Вы правильно выделили $T^4$ как неизвестную величину. При сокращении степеней десяти запись свелась к $T^4=16$, после чего указан ответ $T=2$.

\textbf{Ваш ответ.} $2$.

\Err{} Потерян множитель $10^{12}$ при делении степеней десяти.

\textbf{Где именно.} Начальная формула и подстановка правильные:
\[
8\cdot10^6=(5\cdot10^{-8})\cdot10\cdot T^4.
\]
Первый неправильный результат --- $T^4=16$ без большого степенного множителя.

\textbf{Почему.} Произведение $10^{-8}\cdot10^1=10^{-7}$. При делении $10^6$ на $10^{-7}$ показатели \emph{вычитаются}:
\[
\frac{10^6}{10^{-7}}=10^{6-(-7)}=10^{13}.
\]
Два минуса внутри показателя дают плюс. Простое деление $8$ на $0{,}5$ не учитывает оставшиеся степени десятки.

\Key{} Сначала отделите коэффициенты от степеней десяти. Затем проверьте порядок величины результата.

\textbf{Исправление и полное решение.}
\begin{align*}
8\cdot10^6 &= 5\cdot10^{-8}\cdot10\cdot T^4,\\
8\cdot10^6 &= 5\cdot10^{-7}\cdot T^4,\\
T^4 &= \frac{8\cdot10^6}{5\cdot10^{-7}}
=\frac85\cdot10^{13},\\
T^4 &= 1{,}6\cdot10^{13}
=16\cdot10^{12}
=(2\cdot10^3)^4.
\end{align*}
С учётом $T>0$: $T=2\cdot10^3=2000$.

\textbf{Проверка.}
\[
5\cdot10^{-8}\cdot10\cdot2000^4
=5\cdot10^{-7}\cdot16\cdot10^{12}
=8\cdot10^6.
\]
\Result{} $\boxed{T=2000}$.

\subsection*{Задача на закрепление}
Из формулы $P=\sigma ST^4$ найдите $T>0$, если $P=10^5$, $\sigma=2\cdot10^{-8}$ и $S=5$.
\newpage
\taskhead{10}{Дробная степень и соотношение объёмов}
\textbf{Условие.} $p_1V_1^{5/3}=p_2V_2^{5/3}$. Известно, что $p_2=32p_1$, $V_1=0{,}016$ м$^3$. Найдите $V_2$.

\textbf{Ваше решение.} Правильно подставлено $p_2=32p_1$ и сокращено $p_1$. В следующей строке дробь $0{,}016=16/1000$ возводится в степень $5/3$ с сохранением знаменателя $1000$ без соответствующего возведения в степень. Далее встречается запись $1/2000$, а обведённый ответ --- $0{,}002$.

\textbf{Ваш ответ.} $0{,}002$ м$^3$. Он совпал с правильным, хотя промежуточные вычисления требуют исправления.

\Err{} Степень дроби применена не ко всему числителю и знаменателю.

\textbf{Где именно.} После сокращения $p_1$ верно:
\[
0{,}016^{5/3}=32V_2^{5/3}.
\]
Проблемный переход --- преобразование $0{,}016^{5/3}$ по записи $0{,}016=16/1000$:
\[
\left(\frac{16}{1000}\right)^{5/3}
=\frac{16^{5/3}}{\color{errorred}1000^{5/3}},
\]
поэтому нельзя оставлять знаменатель равным просто $1000$.

\textbf{Почему.} Правило степени дроби:
\[
\left(\frac ab\right)^r=\frac{a^r}{b^r}
\qquad(a>0,\ b>0).
\]
Один и тот же показатель применяется к \emph{обеим} частям дроби. Если показатель применён только к числителю, равенство нарушается.

\Key{} В этой задаче проще вообще не раскрывать десятичную дробь: сначала найдём отношение объёмов.

\textbf{Исправление и полное решение.}
\begin{align*}
p_1V_1^{5/3}&=32p_1V_2^{5/3},\\
V_1^{5/3}&=32V_2^{5/3}\quad(p_1>0),\\
\left(\frac{V_1}{V_2}\right)^{5/3}&=32=2^5,\\
\frac{V_1}{V_2}&=(2^5)^{3/5}=2^3=8,\\
V_2&=\frac{V_1}{8}=\frac{0{,}016}{8}=0{,}002\ \text{м}^3.
\end{align*}
\textbf{Проверка.} Поскольку $V_1/V_2=8$, то
\[
\frac{V_1^{5/3}}{V_2^{5/3}}=8^{5/3}=32=\frac{p_2}{p_1}.
\]
Равенство из условия выполнено.

\Result{} $\boxed{V_2=0{,}002\ \text{м}^3}$.

\subsection*{Задача на закрепление}
В процессе $p_1V_1^{5/3}=p_2V_2^{5/3}$ известно, что $p_2=243p_1$, $V_1=0{,}081$ м$^3$. Найдите $V_2$.
\vfill
\begin{center}\footnotesize\textcolor{textgray}{Лёвин Артём Александрович эксклюзивно для Даниила Хрястова}\end{center}
\end{document}
